\documentclass[11pt]{article}

\usepackage[a4paper,margin=1in]{geometry}
\usepackage{enumitem}
\usepackage{array}
\usepackage{longtable}
\usepackage{booktabs}
\usepackage{hyperref}
\usepackage{xcolor}
\usepackage{titlesec}

\setlist[itemize]{leftmargin=1.8em}
\setlist[enumerate]{leftmargin=2.0em}
\titleformat{\section}{\large\bfseries}{\thesection.}{0.5em}{}
\titleformat{\subsection}{\normalsize\bfseries}{\thesubsection.}{0.5em}{}

\title{Submission Readiness Checklist and 6-Month Execution Plan\\
for Pre-Basis Structural Simplification in Quantum Compilation}
\author{YangJinsey}
\date{\today}

\begin{document}

\maketitle

\section{Project Positioning}

\noindent
This project should be developed along a two-track strategy:

\begin{itemize}
    \item \textbf{Track A: Quantum-ready paper.}  
    The goal is a strong, reproducible compiler-systems paper showing that bounded pre-basis structural simplification can prevent severe regressions on structured workloads and recover strong external-baseline quality on real algorithm-family circuits.
    
    \item \textbf{Track B: PRX-ready upgrade path.}  
    The goal is not merely to improve one codebase, but to elevate the project into a more general principle of quantum compilation: namely, that basis lowering can irreversibly destroy higher-level structure, and that a bounded pre-basis simplification stage can preserve exploitable optimizability gaps for structured circuit families.
\end{itemize}

\section{Master Checklist}

\subsection{A. Must-Do for a Strong Quantum Submission}

\begin{enumerate}
    \item \textbf{Refresh all existing tables on the latest optimized branch.}
    \begin{itemize}
        \item Re-run all headline result tables.
        \item Eliminate stale numbers, timeouts from old code, and inconsistent baselines.
        \item Ensure every figure and table is regenerated from the current implementation.
    \end{itemize}

    \item \textbf{Consolidate the empirical claim.}
    \begin{itemize}
        \item Make the main claim precise:
        \begin{quote}
        bounded pre-basis structural preprocessing can act as an anti-regression and quality-recovery layer for structured circuits in UCC
        \end{quote}
        \item Avoid overclaiming universal superiority.
        \item Explicitly distinguish:
        \begin{itemize}
            \item strong success cases,
            \item anti-regression cases,
            \item weak/no-gain cases.
        \end{itemize}
    \end{itemize}

    \item \textbf{Add at least one public benchmark source beyond current curated Qiskit examples.}
    \begin{itemize}
        \item Preferred: MQT Bench or SupermarQ.
        \item Minimum target: 2--3 representative benchmark instances.
        \item Use the same comparison protocol as the current core tables.
    \end{itemize}

    \item \textbf{Include at least one backend-aware or architecture-aware study.}
    \begin{itemize}
        \item Use a fixed coupling map or a realistic backend target.
        \item Report: total gates, depth, 2-qubit gates, runtime.
        \item Compare: baseline UCC, optimized UCC, Qiskit opt3.
    \end{itemize}

    \item \textbf{Strengthen ablation and runtime credibility.}
    \begin{itemize}
        \item Isolate the role of:
        \begin{itemize}
            \item pre-basis preprocessing,
            \item candidate selection,
            \item short-circuiting,
            \item commutative inverse cancellation,
            \item bounded iteration / bounded block selection.
        \end{itemize}
        \item Add repeated timing for major headline cases if feasible.
    \end{itemize}

    \item \textbf{Provide one clean artifact story.}
    \begin{itemize}
        \item One-command or clearly scripted reproduction.
        \item Explicit mapping from scripts to tables/figures.
        \item Saved JSON/CSV/Markdown outputs.
    \end{itemize}
\end{enumerate}

\subsection{B. High-Value Upgrades for PRX Competitiveness}

\begin{enumerate}
    \item \textbf{Upgrade the central claim from a UCC result to a compiler principle.}
    \begin{itemize}
        \item Move from:
        \begin{quote}
        ``this improves UCC''
        \end{quote}
        to:
        \begin{quote}
        ``there exists a class of structured circuits for which pre-basis simplification is fundamentally more effective than post-lowering simplification''
        \end{quote}
        \item Formalize the notion of structure loss after basis lowering.
    \end{itemize}

    \item \textbf{Broaden compiler comparison beyond UCC and Qiskit alone.}
    \begin{itemize}
        \item Add TKET if feasible.
        \item Add one additional relevant baseline if practical.
        \item Use fair, explicitly documented comparison settings.
    \end{itemize}

    \item \textbf{Expand to a broader public benchmark corpus.}
    \begin{itemize}
        \item Include multiple sources rather than one curated family.
        \item Cover more than three real instances.
        \item Ensure family diversity:
        \begin{itemize}
            \item phase-estimation style,
            \item amplitude-estimation / amplification style,
            \item QAOA beyond a single ansatz family,
            \item mirrored / inverse-structured search circuits.
        \end{itemize}
    \end{itemize}

    \item \textbf{Show repeated wins, not just parity.}
    \begin{itemize}
        \item A PRX-style paper should ideally contain multiple real families where your method:
        \begin{itemize}
            \item beats strong baselines in gate count or depth, or
            \item beats them in runtime at matched quality, or
            \item yields a hardware-relevant 2Q advantage.
        \end{itemize}
    \end{itemize}

    \item \textbf{Add a stronger theoretical section.}
    \begin{itemize}
        \item Define classes of circuits with pre-lowering vs post-lowering optimizability gaps.
        \item State sufficient conditions for recoverable vs irrecoverable structure loss.
        \item Analyze bounded preprocessing complexity and the placement gap conceptually.
    \end{itemize}

    \item \textbf{Tie empirical results to a general scientific conclusion.}
    \begin{itemize}
        \item The paper should conclude more than:
        \begin{quote}
        ``this branch works well''
        \end{quote}
        \item It should conclude something like:
        \begin{quote}
        ``compiler pass placement is itself a scientifically meaningful variable in structured quantum circuit optimization''
        \end{quote}
    \end{itemize}
\end{enumerate}

\subsection{C. Things Not Worth Spending Too Much Time On}

\begin{enumerate}
    \item polishing already-solved toy cases;
    \item adding many more synthetic variants without expanding scientific insight;
    \item chasing tiny constant-factor speedups before the paper-level claim is stable;
    \item over-engineering code cleanup before the final experimental story is fixed;
    \item spending weeks on prose before the missing experiments and theory framing are done.
\end{enumerate}

\section{Priority Triage}

\subsection{Tier 1: Must Finish}
\begin{itemize}
    \item refresh all main tables on latest code;
    \item finalize ablation refresh;
    \item add one public benchmark source;
    \item add one backend-aware study;
    \item stabilize artifact reproduction;
    \item write clean problem statement and main claim.
\end{itemize}

\subsection{Tier 2: Strongly Recommended}
\begin{itemize}
    \item repeated timing / variance on headline cases;
    \item one additional real algorithm family beyond the current three;
    \item failure-case section;
    \item a stronger introduction with a principled statement about structure loss.
\end{itemize}

\subsection{Tier 3: PRX-Oriented Stretch Goals}
\begin{itemize}
    \item cross-compiler comparison beyond Qiskit;
    \item formal statement of pre-basis advantage conditions;
    \item multiple public benchmark suites;
    \item repeated clear wins on real workloads;
    \item a theory-driven section on placement and optimizability gaps.
\end{itemize}

\section{Six-Month Detailed Schedule}

\noindent
The schedule below is written at a two-week granularity. It is designed so that:
\begin{itemize}
    \item by Month 3, the project is in strong \textbf{Quantum-ready} shape;
    \item by Month 6, the project is either:
    \begin{itemize}
        \item submitted to \textbf{Quantum} / \textbf{ACM TQC} / \textbf{IEEE TQE}, or
        \item upgraded as far as realistically possible toward \textbf{PRX Quantum}.
    \end{itemize}
\end{itemize}

\begin{longtable}{p{0.10\textwidth} p{0.20\textwidth} p{0.32\textwidth} p{0.30\textwidth}}
\toprule
\textbf{Weeks} & \textbf{Theme} & \textbf{Main Tasks} & \textbf{Deliverables} \\
\midrule
\endhead

1--2 &
Audit and stabilization &
\begin{itemize}
    \item audit all existing result files;
    \item identify stale tables and mismatched branches;
    \item freeze the current optimized branch for paper experiments;
    \item define the exact baseline protocol.
\end{itemize}
&
\begin{itemize}
    \item experiment inventory;
    \item list of tables to rerun;
    \item frozen experiment configuration file.
\end{itemize}
\\

3--4 &
Core refresh &
\begin{itemize}
    \item rerun headline real-instance tables;
    \item rerun scaling tables;
    \item rerun ablation tables;
    \item verify reproducibility of each table.
\end{itemize}
&
\begin{itemize}
    \item refreshed master result tables;
    \item regenerated figures;
    \item updated results directory with current outputs.
\end{itemize}
\\

5--6 &
Public benchmark expansion I &
\begin{itemize}
    \item integrate one public benchmark source;
    \item select 2--3 representative instances;
    \item run baseline UCC / optimized UCC / Qiskit opt3;
    \item inspect whether current claim generalizes.
\end{itemize}
&
\begin{itemize}
    \item one new benchmark subsection;
    \item raw data for public benchmark results.
\end{itemize}
\\

7--8 &
Hardware-aware study &
\begin{itemize}
    \item choose one realistic coupling map or backend target;
    \item define routing / compilation settings fairly;
    \item run hardware-aware comparisons;
    \item record depth and 2Q metrics.
\end{itemize}
&
\begin{itemize}
    \item one backend-aware table;
    \item one backend-aware discussion subsection.
\end{itemize}
\\

9--10 &
Ablation deepening and timing quality &
\begin{itemize}
    \item isolate method components more clearly;
    \item add repeated timing for headline cases;
    \item summarize no-gain / weak-gain cases;
    \item identify runtime bottlenecks in the current implementation.
\end{itemize}
&
\begin{itemize}
    \item cleaned ablation section;
    \item timing variance table;
    \item shortlist of implementation bottlenecks.
\end{itemize}
\\

11--12 &
Paper draft I: Quantum version &
\begin{itemize}
    \item write abstract, introduction, related work;
    \item write method section;
    \item write experiment protocol and main results;
    \item draft limitations and failure cases.
\end{itemize}
&
\begin{itemize}
    \item complete first full draft for a Quantum-style paper.
\end{itemize}
\\

13--14 &
Narrative upgrade &
\begin{itemize}
    \item rewrite the introduction around the structure-loss problem;
    \item define placement as a central variable;
    \item sharpen the distinction between:
    \begin{itemize}
        \item anti-regression,
        \item quality recovery,
        \item true external-baseline wins.
    \end{itemize}
\end{itemize}
&
\begin{itemize}
    \item stronger framing memo;
    \item revised abstract and introduction.
\end{itemize}
\\

15--16 &
Theory framing I &
\begin{itemize}
    \item formalize circuit-family categories:
    \begin{itemize}
        \item exact inverse,
        \item repeated exact,
        \item mirrored,
        \item commuting inverse chains.
    \end{itemize}
    \item write a clean complexity discussion;
    \item define recoverable vs irrecoverable structure loss.
\end{itemize}
&
\begin{itemize}
    \item theory note or appendix draft;
    \item candidate theorem/proposition statements.
\end{itemize}
\\

17--18 &
PRX-oriented expansion I &
\begin{itemize}
    \item test one additional real algorithm family;
    \item try one more compiler baseline if feasible;
    \item examine whether any result becomes a genuine repeated win.
\end{itemize}
&
\begin{itemize}
    \item expanded comparison table;
    \item decision memo: PRX upgrade still plausible or not.
\end{itemize}
\\

19--20 &
Implementation optimization pass &
\begin{itemize}
    \item optimize the worst runtime bottlenecks only if they affect headline claims;
    \item improve bounded scanning, block selection, or caching where most valuable;
    \item rerun the few cases most affected by speed improvements.
\end{itemize}
&
\begin{itemize}
    \item one optimized implementation revision;
    \item before/after timing comparison.
\end{itemize}
\\

21--22 &
Paper draft II: submission version &
\begin{itemize}
    \item integrate all refreshed results;
    \item finalize figures and captions;
    \item finalize artifact appendix;
    \item decide target venue based on evidence.
\end{itemize}
&
\begin{itemize}
    \item near-submission manuscript;
    \item complete artifact guide.
\end{itemize}
\\

23--24 &
External review and final polish &
\begin{itemize}
    \item send draft to trusted readers / mentors;
    \item incorporate feedback;
    \item reduce overclaiming;
    \item finalize author list, acknowledgments, and repository packaging.
\end{itemize}
&
\begin{itemize}
    \item final submission package;
    \item arXiv-ready source;
    \item journal submission decision.
\end{itemize}
\\

\bottomrule
\end{longtable}

\section{Monthly Milestones}

\subsection*{End of Month 1}
\begin{itemize}
    \item all stale experiments identified;
    \item core tables refreshed;
    \item experiment pipeline stabilized.
\end{itemize}

\subsection*{End of Month 2}
\begin{itemize}
    \item one public benchmark source added;
    \item one hardware-aware study completed;
    \item strong Quantum-level empirical story assembled.
\end{itemize}

\subsection*{End of Month 3}
\begin{itemize}
    \item full Quantum-style draft completed;
    \item artifact path clear;
    \item realistic submission to Quantum / ACM TQC / IEEE TQE becomes possible.
\end{itemize}

\subsection*{End of Month 4}
\begin{itemize}
    \item framing upgraded from ``UCC optimization'' to ``structure-loss and placement principle'';
    \item theory note drafted.
\end{itemize}

\subsection*{End of Month 5}
\begin{itemize}
    \item extra compiler or benchmark evidence added if feasible;
    \item runtime bottlenecks selectively improved;
    \item PRX viability reassessed honestly.
\end{itemize}

\subsection*{End of Month 6}
\begin{itemize}
    \item either:
    \begin{itemize}
        \item submit the strongest realistic version to Quantum / ACM TQC / IEEE TQE,
    \end{itemize}
    \item or:
    \begin{itemize}
        \item if the generality, theory, and repeated wins are now strong enough, prepare a PRX Quantum submission.
    \end{itemize}
\end{itemize}

\section{Decision Rule at the End of the Plan}

\noindent
At the end of the six-month cycle, choose the venue using the following rule:

\begin{itemize}
    \item \textbf{Submit to Quantum / ACM TQC / IEEE TQE} if the paper has:
    \begin{itemize}
        \item refreshed and credible experiments,
        \item one public benchmark extension,
        \item one backend-aware study,
        \item a strong reproducibility package,
        \item a clean compiler-systems claim,
        \item but mostly parity / anti-regression rather than repeated wins.
    \end{itemize}

    \item \textbf{Consider PRX Quantum} only if the paper has:
    \begin{itemize}
        \item a clear general principle beyond UCC,
        \item broad public evidence across multiple families,
        \item strong external-baseline wins on real workloads,
        \item hardware-relevant validation,
        \item and a meaningful theoretical component.
    \end{itemize}
\end{itemize}

\end{document}